\documentclass[11pt]{article}
\usepackage[margin=27mm]{geometry}
\usepackage[T1]{fontenc}
\usepackage{lmodern}
\usepackage{amsmath,amssymb,amsthm,booktabs,microtype}
\usepackage[colorlinks=true,allcolors=blue]{hyperref}
\newtheorem{theorem}{Theorem}
\newtheorem{lemma}[theorem]{Lemma}
\newtheorem{corollary}[theorem]{Corollary}
\newcommand{\QQ}{\mathbb Q}
\newcommand{\CC}{\mathbb C}
\newcommand{\FF}{\mathbb F}
\newcommand{\ZZ}{\mathbb Z}
\newcommand{\ff}[2]{#1^{\underline{#2}}}
\setlength{\emergencystretch}{2em}
\title{Ancillary proof: continuous-convolution rigidity through degree 72}
\author{Alexander Migita}
\date{}
\begin{document}
\maketitle

For \(A\in\CC[x]\), put
\[
 H_A(x)=\int_0^xA(t)A(x-t)\,dt
 =\sum_{i,j}a_ia_j\frac{i!j!}{(i+j+1)!}x^{i+j+1},
 \qquad A=\sum_i a_ix^i.
\]
The integral is a polynomial identity.
We prove the finite degree range used in the manuscript.
The accompanying JSON supplies integer polynomial identities; the
standalone replay uses exact arithmetic and checks every degree and
every residual stratum. No assertion of rigidity in all degrees, or of
simplicity of the discrete matching system, follows from this note.

\begin{theorem}
Every nonzero complex polynomial of degree at most 72 satisfying
\(A\mid H_A\) is a monomial.
\end{theorem}

\section{The uniform prime-power reduction}

Let \(\ell\) be an odd prime, \(e\ge1\), \(P=\ell^e\), and
\[
 2d+1=P+g,\qquad 0\le g<\ell,\qquad g\text{ even}.
\]
Write \(A=x^d+a_1x^{d-1}+\cdots+a_d\), with \(a_0=1\).
For \(g>0\), put \(\delta=(g-1)/2\) and
\[
 S_{g,A}(x)=
 \sum_{r+s\le g}
 \frac{\ff g{r+s}}{\ff\delta r\,\ff\delta s}
 a_ra_sx^{g-r-s}.
\]
For \(g=0\), set \(S_{0,A}=1\).

\begin{lemma}
The polynomial \(P H_A\) has coefficients integral at \(\ell\), and
\[
 P H_A(x)\equiv u x^P S_{g,A}(x)\pmod\ell,\qquad
 u=\overline{\frac{P(d!)^2}{(P+g)!}}\ne0.
\]
The assertion is an identity in the coefficient indeterminates of \(A\).
\end{lemma}

\begin{proof}
For \(\beta(i,j)=i!j!/(i+j+1)!\), Legendre's formula gives
\[
 v_\ell(\beta(i,j)^{-1})=
 \sum_{h\ge1}
 \left(\left\lfloor\frac{i+j+1}{\ell^h}\right\rfloor
       -\left\lfloor\frac i{\ell^h}\right\rfloor
       -\left\lfloor\frac j{\ell^h}\right\rfloor\right).
\]
Each summand is zero or one.
Since \(i+j+1\le P+g<\ell P\), all summands with \(h>e\) vanish.
This proves integrality after multiplication by \(P\).
If \(i+j+1<P\), the summand at \(h=e\) also vanishes, so the
normalized coefficient reduces to zero.

For a remaining term, write \(i=d-r,\ j=d-s\), \(r+s\le g\).
For \(1\le h\le e\),
\[
 d\bmod\ell^h=\frac{\ell^h+g-1}{2}.
\]
Subtracting \(r,s\le g<\ell\) creates no borrow, and the sum of the
two residues plus one is \(\ell^h+g-r-s\ge\ell^h\).
All \(e\) summands are therefore one.
In particular the leading normalized coefficient is a unit.
Dividing by that leading coefficient gives
\[
 \frac{\ff{(P+g)}{r+s}}{\ff d r\,\ff d s}
 \equiv
 \frac{\ff g{r+s}}{\ff\delta r\,\ff\delta s}\pmod\ell.
\]
The denominators are units because their odd numerators have absolute
value less than \(\ell\). This proves the coefficient formula.
\end{proof}

\begin{lemma}
Suppose that for each \(1\le k\le g\) no monic degree-\(k\)
polynomial \(C\), with \(C(0)\ne0\), satisfies
\[
 C\mid S_{g,C}
 \quad\text{over }\overline{\FF}_\ell.
\]
Then rigidity holds in degree \(d\).
For \(g=0\), no residual hypothesis is needed.
\end{lemma}

\begin{proof}
A hypothetical nonmonomial solution gives an algebraic solution:
the remainder equations have rational coefficients, and adjoining
the inverse of a nonzero coefficient produces an affine system with
a point over \(\overline{\QQ}\).
At a place above \(\ell\), give the descending coefficient \(a_j\)
weight \(j\), and put
\(\eta=\min_{a_j\ne0}v_\ell(a_j)/j\).
After a finite extension choose \(t\) with \(v_\ell(t)=-\eta\).
Then \(t^dA(x/t)\) is monic and integral, with a nonleading unit.
This scaling preserves divisibility, since its convolution is
\(t^{2d+1}H_A(x/t)\).

Monic division of the integral normalized convolution has integral
quotient. The preceding lemma therefore implies
\(\overline A\mid x^P S_{g,\overline A}\).
If \(g=0\), this forces \(\overline A=x^d\), contradicting the
nonleading unit.
Otherwise write
\(\overline A=x^{d-k}C\), \(C(0)\ne0\).
The degree-\(g\) residual forces \(1\le k\le g\).
Its descending coefficient list is the one used for \(C\), so
\(C\mid S_{g,C}\), contrary to the hypothesis.
The argument permits arbitrary ramification.
\end{proof}

\section{The finite certificates}

For \(g=2,4,6\) and every \(1\le k\le g\), normalize \(C(0)=1\):
\[
 C=x^k+c_1x^{k-1}+\cdots+c_{k-1}x+1.
\]
This loses no geometric point over the algebraic closure.
Indeed \(C_t(x)=t^kC(x/t)\) preserves the residual divisibility,
and a suitable \(k\)-th root makes its constant coefficient one.

Take \emph{all} coefficients of
\(\operatorname{rem}_C S_{g,C}\), and clear their rational
denominators, without discarding integer content.
Call the resulting integer polynomials \(F_{g,k,j}\).
The file \texttt{certificate.json} contains integer polynomials
\(U_{g,k,j}\) and positive integers \(D_{g,k}\) with
\[
 \sum_j U_{g,k,j}F_{g,k,j}=D_{g,k}.
\]
When \(\ell\) divides neither \(D_{g,k}\) nor the original
denominator scales, this identity rules out common zeros over
every extension of \(\FF_\ell\).
The JSON also contains the complete prime factorizations of these
integers and recursive Lucas primality certificates.

For example,
\[
 S_{2,C}=x^2+8c_1x+8c_1^2-16c_2.
\]
In degree one the remainder is \(c_1^2\); in degree two it is
\(7c_1x+8c_1^2-17c_2\).
The only possible exceptional odd primes greater than two in this
criterion are 7 and 17.
The same certificate procedure at gaps four and six supplies finite
exceptional-prime covers.
An exceptional prime is not declared a counterexample to
characteristic-zero rigidity.

The stored coverage table gives, for every \(1\le d\le72\),
a triple \((\ell,e,g)\), \(g\in\{0,2,4,6\}\), such that
\[
 2d+1=\ell^e+g,\qquad \ell>g,
\]
and, for \(g>0\), \(\ell\) avoids every relevant certificate factor
in every degree \(1\le k\le g\).
For illustration, three choices filling gaps between the
prime-index cases are
\[
 95=89+6,\qquad119=113+6,\qquad123=11^2+2.
\]
The theorem follows from the full 72-row table and the preceding
lemmas. Degree zero is immediate.

\section{Independent replay}

Run
\begin{verbatim}
python verify_coverage.py
\end{verbatim}
in this directory, with SymPy installed.
The program reconstructs the residual equations from the displayed
falling-factorial kernel, checks every integer identity and
denominator scale, verifies the prime factorizations and Lucas
certificates, and validates every row of the coverage table.
As an additional direct check, it rebuilds all 132348 ordered
bilinear coefficients of \(P H_A\) for the 72 chosen degrees and
compares them with the lemma.
No computer algebra system beyond SymPy is required.

The final report explicitly records that the conclusion concerns
\(0\le\deg A\le72\), corresponding to ODE orders \(2\le p\le73\)
in the manuscript. It does not assert simplicity of those discrete
matching schemes. The main paper's prime-index simplicity theorem
and its separate order-five calculation supply their own arguments.
\end{document}
